\documentclass[11pt]{article}
\usepackage[a4paper,margin=25mm]{geometry}
\usepackage{amsmath,amssymb}
\usepackage[T1]{fontenc}
\usepackage{lmodern}
\usepackage{hyperref}
\hypersetup{colorlinks=true,urlcolor=blue}
\setlength{\parindent}{0pt}
\setlength{\parskip}{7pt}
\title{A complete mapping witness for the conjecture as written}
\author{SucRunBug}
\date{1 October 2026}
\begin{document}
\maketitle
\textbf{Conjecture 00000000127 -- Verdict: TRUE AS WRITTEN.}

\textbf{Filed statement.} For every prime $p>3$, some integer exponent $d\notin\{1,p-2\}$ makes both $x\mapsto x^d$ and $x\mapsto x^d+x$ permutations of $\mathbb F_p$.

\section*{Proof}
Choose $d=p$. Since $p>3$, this is a positive integer distinct from both $1$ and $p-2$. Fermat's identity on the prime field states that $x^p=x$ for every $x\in\mathbb F_p$, including $x=0$. The two maps are therefore
\[x\longmapsto x,\qquad x\longmapsto 2x.\]
The first is the identity permutation. The second is a permutation because $p>3$ implies $2\ne0$ in the field; its inverse is multiplication by $2^{-1}$. This proves exactly the quantified statement in the file.

\paragraph{Interpretation boundary.} The filed conjecture does not impose $d<p$, $1\le d\le p-2$, or a requirement that $d$ be a canonical representative modulo $p-1$. Our choice is equivalent as a polynomial \emph{function} to exponent $1$, which is why it provides no nontrivial new complete mapping. If one adds a requirement excluding exponents congruent to $1$ modulo $p-1$, or restricts the allowed exponent range, this witness is excluded and the present proof does not settle that modified question. We state this prominently so that the submission cannot be mistaken for a result on nontrivial complete mappings.

\section*{Lean formalization}
The Lean theorem uses the standard prime field \texttt{ZMod p}, quantifies over every prime $p>3$, explicitly proves both exclusions on $d$, and proves bijectivity of both functions. The exponent type is $\mathbb N$, reflecting the literal unrestricted-exponent reading.

\textbf{Reproduction.} In the supplied \texttt{lean4/} project, run
\texttt{lake exe cache get}, then \texttt{lake build} and
\texttt{lake env lean Check.lean}. The pinned versions are Lean and Mathlib
\texttt{v4.33.0}. The supplied \texttt{reproduce.py} performs independent
finite sanity checks; these checks are not the proof of any infinite assertion.

\textbf{Source.} \url{https://github.com/The-Last-Math-Competition/The-Last-Math-Competition/blob/main/conjectures/00000000127.md}
\end{document}
